%=====================================================================
% 10 — KẾT LUẬN — scientific-paper compression
%=====================================================================
\section{Kết luận}
\label{sec:conclusion}

Bài báo đề xuất một kiến trúc tham chiếu phần mềm tỉnh--trung ương cho nền tảng quản lý DN KH\&CN và DN KNST, trong đó quyền xử lý nghiệp vụ chủ yếu thuộc cấp tỉnh nhưng dữ liệu, liên thông và tổng hợp phải bảo đảm tính thống nhất ở cấp Bộ/quốc gia. Đóng góp chính là một mô hình \emph{gán quyền ghi theo phạm vi thẩm quyền và thời gian hiệu lực} -- với máy trạng thái chuyển giao và bất biến single-active-writer -- tổng quát hóa được cho các nền tảng chính phủ đa cấp khác, không riêng miền DN KH\&CN/DN KNST; cách tổ chức \emph{lõi ổn định} và \emph{điểm biến thiên} của SRA; ranh giới trách nhiệm liên thông; và chuỗi truy vết cho phép nối căn cứ nguồn với quyết định kiến trúc và tiêu chí đánh giá. SRA không bắt buộc một topology lưu trữ; P-DIST chỉ được dùng như một cấu hình để thử thách các bất biến kiến trúc dưới điều kiện phân tán. Ba nguyên tắc thiết kế được rút ra ở mức khái quát là: phân rã thẩm quyền theo nhóm thuộc tính--phạm vi--thời gian; tách lõi pháp lý ổn định khỏi biến thể vận hành có kiểm soát; và duy trì provenance/liên thông độc lập với topology lưu trữ.

Mười kịch bản cho thấy ứng viên ban đầu còn hai khoảng trống quan trọng: chuyển quyền ghi phải được biểu diễn theo thời gian và kiến trúc phải chịu được giai đoạn các đơn vị vận hành khác phiên bản. Sau khi hai vấn đề này được hấp thụ vào B1-R, hậu kiểm bằng cùng rubric cho kết quả 2 kịch bản đáp ứng trực tiếp, 8 kịch bản đáp ứng có điều kiện và không còn rủi ro kiến trúc chưa được hấp thụ ở mức L2 trong bộ kịch bản hiện tại. Kết quả này là \emph{scenario-based architecture verification} trước triển khai, không phải bằng chứng thực nghiệm về một hệ thống vận hành. Bước tiếp theo là xây dựng nguyên mẫu và kiểm thử lặp lại đối với chuyển quyền ghi, nâng cấp lệch phiên bản, liên thông, phục hồi, an toàn và các thuộc tính chất lượng đo được.
